% chapters/ch14_python_automation.tex — Chapter 14: Python for Test Automation
% Scope (PLAN.md): struct.pack/unpack for binary telemetry, pyserial scripting,
% log parsing & regex, CSV/report generation, the questions asked when "Python
% automation framework" is on a resume, 15 graded problems.
% Budget 25 pp. All pure-logic Python verified with python3 (runs clean);
% pyserial/hardware snippets are marked as requiring the module/target.
\chapter{Python for Test Automation}
\label{ch:python-automation}

\begin{partnote}
\textbf{Scope} --- \code{struct.pack}/\code{unpack} of binary telemetry
$\rightarrow$ \code{pyserial} port scripting $\rightarrow$ log-file parsing \&
regex $\rightarrow$ CSV / report generation $\rightarrow$ the questions you
\emph{will} be asked when ``Python automation'' is on your resume --- with 15
graded problems. This is straight off the resume: the Python suite parsing
MIL-STD-1553B and RS422 telemetry that cut regression effort by ~60\%. Expect
the panel to make you defend exactly how that worked.
\end{partnote}

When ``Python automation framework'' appears on an embedded resume, the
interviewer's job is to find out whether you \emph{built} it or just ran it.
The questions probe binary protocol handling (the hard part --- bytes, not
strings), serial I/O, parsing real device logs, and turning results into a
report --- and whether your framework was maintainable. Every problem here is a
piece of that.

%=====================================================================
\section{Primer}
%=====================================================================

\subsection{Binary data: \code{struct} is the whole game}

Test automation for embedded means handling \emph{bytes}, and Python's
\code{struct} module packs/unpacks C-style binary data to/from \code{bytes}
objects with a format string: \code{>BBfH} is big-endian (\code{>}) then an
unsigned byte, byte, 32-bit float, 16-bit unsigned --- exactly a telemetry
frame. The endianness prefix (\code{>} big / \code{<} little / \code{!}
network) is mandatory and the \#1 source of bugs (Chapters 1/8): get it wrong
and every multi-byte field is byte-swapped garbage. Sub-byte fields (a 5-bit RT
address) \code{struct} can't do directly --- you unpack the containing word then
extract with shifts and masks, the same discipline as the C side.

\subsection{Serial I/O with \code{pyserial}}

\code{pyserial} (\code{import serial}) opens a port
(\code{serial.Serial('/dev/ttyUSB0', 115200, timeout=1)}) and gives
\code{read}/\code{write}/\code{read\_until}. The automation pattern: write a
command frame (packed with \code{struct}), read the response with a
\emph{timeout} (never block forever on hardware), parse it, and assert. Timeouts
and framing (find the sync, read the declared length) matter as much as on the
firmware side --- a test harness that hangs on a missing byte is useless.

\subsection{Parsing logs and generating reports}

Real device output is messy text; \textbf{regular expressions}
(\code{import re}) extract the structured bits --- timestamps, error codes,
measured values --- from log lines. Results then become artefacts: a
\textbf{CSV} (\code{import csv}) for data analysis or import into a traceability
tool, and a human-readable \textbf{report} (text/markdown/HTML) summarising
pass/fail. The framework's value is in this loop: \emph{drive the hardware
$\rightarrow$ capture $\rightarrow$ parse $\rightarrow$ compare to expected
$\rightarrow$ report}, run automatically and repeatably --- which is how a manual
regression that took days becomes a script that takes minutes (the resume's
60\%).

%=====================================================================
\section{Q\&A Bank}
%=====================================================================

\tierbanner{Tier 1 --- Screening}{Two to four sentences.}

\question{Q1.1}{How do you handle binary protocol data in Python?}
With the \code{struct} module: \code{struct.pack(fmt, ...)} builds a
\code{bytes} frame and \code{struct.unpack(fmt, data)} parses one, where the
format string encodes endianness and each field's type. Sub-byte fields are
extracted with shifts and masks after unpacking the containing integer.

\question{Q1.2}{Why does endianness matter in \code{struct} formats?}
Because the format's leading character (\code{>}, \code{<}, \code{!}, \code{=})
sets byte order, and a wrong choice byte-swaps every multi-byte field ---
turning a valid telemetry value into garbage. Avionics buses are often
big-endian while the host is little-endian, so you specify it explicitly.

\question{Q1.3}{What library do you use for serial port automation?}
\code{pyserial} --- \code{serial.Serial()} (with a \code{timeout}) to open, then
\code{write()} to send and \code{read()} / \code{read\_until()} to receive. A timeout is
essential so the test never blocks forever on missing hardware data.

\question{Q1.4}{How do you extract fields from a messy log line?}
With a regular expression (\code{re} module) that captures the structured parts
--- timestamp, level, code, value. \code{re.match}/\code{re.search} with capture
groups pulls them out for comparison against expected values.

\question{Q1.5}{How do you report automation results?}
Aggregate pass/fail counts into a summary, write the raw data to CSV for
analysis/traceability, and produce a human-readable report (text/markdown/HTML).
The report is the artefact that makes the run auditable.

\tierbanner{Tier 2 --- Core technical}{A short paragraph.}

\question{Q2.1}{Walk through parsing a binary telemetry frame end to end.}
Read the bytes (from serial or a capture file); find the frame by its sync
pattern and read the declared length (don't trust alignment); \code{struct.unpack}
the fixed fields with the correct endianness; extract any sub-byte fields with
shifts/masks; \emph{validate} the CRC/checksum over the frame and the status/SSM
before using the values; then compare against expected with a tolerance and
record pass/fail. It's the exact pipeline of the C parser (Chapter 8), in Python
--- which is why understanding the protocol matters more than the language.

\question{Q2.2}{Your \code{struct.unpack} gives nonsense values --- debug
checklist?}
Check, in order: \textbf{endianness} (wrong prefix byte-swaps everything ---
the most common cause); \textbf{format string} matches the actual field
sizes/types and their order; \textbf{alignment/padding} (use an explicit
\code{>}/\code{<}/\code{!} prefix, which also disables Python's native
\emph{padding} --- a frequent surprise); the \textbf{length} of the bytes object
matches \code{struct.calcsize(fmt)}; and that you found the true frame start
(sync), not mid-frame. Print a hexdump of the raw bytes and decode by hand for
the first frame to anchor it.

\question{Q2.3}{How do you make a serial test robust against missing/garbled
data?}
Always set a \code{timeout} so reads return rather than block; frame on the sync
word and the length field, not on byte count alone; validate the CRC and
\emph{re-synchronise} on failure (scan for the next sync) instead of trusting
alignment; and treat a timeout as a defined test failure, not a hang. Wrap the
port in try/finally so it's closed even on exception. This mirrors the firmware
receiver's discipline (Chapter 5).

\question{Q2.4}{How did automation cut regression effort ~60\% --- what's the
mechanism?}
Manual regression means a human sending stimuli and eyeballing results for every
test, every build --- slow, error-prone, and not repeatable. The framework
encodes each test as: drive the input (packed command over serial), capture the
response, parse and compare to the expected value automatically, and log
pass/fail. Running the full suite then costs minutes of unattended execution
instead of days of manual work, and it's \emph{repeatable and traceable} (the
same inputs, recorded results) --- which is the actual saving and why it also
improves quality, not just speed.

\question{Q2.5}{How do you keep an automation framework maintainable?}
Separate concerns: a protocol layer (pack/unpack, CRC), a transport layer
(serial/file I/O), the test cases (data-driven --- each test a row of
input/expected, not copy-pasted code), and reporting. Drive tests from data
(CSV/JSON tables) so adding a case is adding a row; use a real test runner
(\code{pytest}/\code{unittest}) for fixtures and reporting; and version the
expected results. The anti-pattern is one giant script with hard-coded frames ---
unmaintainable the moment the protocol changes.

\tierbanner{Tier 3 --- Trap \& depth}{A paragraph plus the trap.}

\question{Q3.1}{Your Python CRC doesn't match the firmware's CRC on the same
bytes. Where do you look?}
CRC mismatches between two correct-looking implementations are almost always a
\textbf{parameter} difference, not a code bug: the \emph{polynomial}, the
\emph{initial value} (0x0000 vs 0xFFFF), whether input/output bits are
\emph{reflected} (bit-reversed), a final \emph{XOR}, and \emph{which bytes} are
covered (does the firmware include the header/length, or start after the sync?).
Also byte order of a multi-byte CRC field on the wire. \trap the trap is staring
at the algorithm; the disciplined move is to pin all the CRC parameters (poly,
init, refin, refout, xorout) on both sides against a known test vector ---
``\code{123456789}'' has published CRC values per variant --- and confirm the two
implementations cover the identical byte range. Matching a Python CRC to the C
\code{crc16} of Chapter 13 (both 0x1021, init 0xFFFF) is exactly this exercise,
and it's the cross-check the resume's parser relies on.

\question{Q3.2}{Is Python fast/deterministic enough for hardware test, and where
are its limits?}
Python is ideal for the \emph{test host} side --- orchestrating, parsing,
comparing, reporting --- where development speed and rich libraries matter and
millisecond-level timing is fine. It is \emph{not} suited to hard real-time
stimulus/response timing (the GC and interpreter add jitter), so cycle-accurate
or sub-millisecond-deterministic test steps belong on the target, an FPGA, or a
real-time HIL rig, with Python commanding them. \trap the trap is proposing
Python to \emph{generate} a precisely-timed waveform or to meet a real-time
deadline; the senior answer puts the determinism in hardware/firmware and uses
Python as the flexible \emph{controller} and \emph{analyser} around it --- which
is exactly its role in a V\&V toolchain (Chapter 12) and on the resume.

\question{Q3.3}{An interviewer says ``you have Python automation on your resume
--- prove you built it, not just ran it.'' What demonstrates depth?}
Concrete design decisions, not features: how you handled \emph{framing} and
\emph{resync} on a noisy serial link; how you matched the firmware's exact CRC
parameters (Q3.1); how you made it \emph{data-driven} so new tests are table
rows; how you handled \emph{timeouts and partial frames} so the suite never
hangs; how you validated the \emph{SSM/status} before trusting telemetry values;
and the measured outcome (regression time before/after). \trap the trap is
listing libraries (``I used pyserial and struct''); depth is in the \emph{design
choices and the edge cases you handled} --- the byte-level framing, the
CRC-parameter matching, the maintainability structure. Being able to whiteboard
the unpack-validate-compare pipeline and name the traps \emph{is} the proof, and
it's the honest, defensible version of the resume claim.

%=====================================================================
\section{Coding Gym --- 15 Python Problems}
%=====================================================================

\begin{partnote}
The pure-logic problems below were run under \code{python3} and their outputs
confirmed; the \code{pyserial} problem is marked as requiring the module and a
target. The CRC-16 here matches the C \code{crc16} of Chapter 13 byte-for-byte
(\code{0x5BCE} on \code{"123"}).
\end{partnote}

%---------------------------------------------------------------
\begin{gymproblem}{Unpack a binary telemetry frame}
Frame: \code{id} (u8), \code{seq} (u8), \code{value} (float), \code{crc} (u16),
big-endian.
\begin{pygym}
import struct

def unpack_frame(b):
    fid, seq, val, crc = struct.unpack('>BBfH', b)   # > = big-endian
    return fid, seq, val, crc
\end{pygym}
\textbf{Note.} The \code{>} prefix is mandatory --- avionics frames are usually
big-endian while the host is little-endian. \code{struct.calcsize('>BBfH')} is 8.
\end{gymproblem}

%---------------------------------------------------------------
\begin{gymproblem}{Pack a command frame}
\begin{pygym}
def pack_cmd(fid, seq, value):
    return struct.pack('>BBf', fid, seq, value)
# pack_cmd(0x10, 1, 3.5) -> b'\x10\x01\x40\x60\x00\x00'
\end{pygym}
\textbf{Note.} The inverse of unpack; this is the stimulus you write to the
serial port to drive a test.
\end{gymproblem}

%---------------------------------------------------------------
\begin{gymproblem}{Extract sub-byte fields (1553 command word)}
\begin{pygym}
def decode_cmd1553(w):
    return dict(rt=(w >> 11) & 0x1F, tr=(w >> 10) & 1,
                sa=(w >> 5) & 0x1F, wc=w & 0x1F)
# decode_cmd1553(0x2C22) -> {'rt': 5, 'tr': 1, 'sa': 1, 'wc': 2}
\end{pygym}
\textbf{Note.} \code{struct} unpacks the 16-bit word; shifts/masks pull the
sub-byte fields --- the same field extraction as the C decoder (Chapter 8).
\end{gymproblem}

%---------------------------------------------------------------
\begin{gymproblem}{CRC-16-CCITT (match the firmware)}
\begin{pygym}
def crc16(data):
    c = 0xFFFF
    for byte in data:
        c ^= byte << 8
        for _ in range(8):
            c = ((c << 1) ^ 0x1021) & 0xFFFF if (c & 0x8000) else (c << 1) & 0xFFFF
    return c
# crc16(b'\x31\x32\x33') -> 0x5BCE  (identical to the C version, Chapter 13)
\end{pygym}
\textbf{Note.} Same poly (0x1021) and init (0xFFFF) as the C \code{crc16} ---
matching CRC parameters across host and target is the cross-check that validates
the parser (Q3.1).
\end{gymproblem}

%---------------------------------------------------------------
\begin{gymproblem}{Parse a log line with regex}
\begin{pygym}
import re

def parse_log_line(line):
    m = re.match(r'\[(\d{2}:\d{2}:\d{2})\]\s+(\w+)\s+code=(\d+)', line)
    return (m.group(1), m.group(2), int(m.group(3))) if m else None
# "[12:34:56] ERROR code=42 ..." -> ('12:34:56', 'ERROR', 42)
\end{pygym}
\textbf{Note.} Capture groups pull timestamp/level/code from messy device output;
\code{None} on no match lets the caller skip non-matching lines.
\end{gymproblem}

%---------------------------------------------------------------
\begin{gymproblem}{Summarise pass/fail}
\begin{pygym}
def summarize(results):
    p = sum(1 for r in results if r)
    return dict(total=len(results), passed=p, failed=len(results) - p,
                rate=round(100.0 * p / len(results), 1))
# [True, True, False, True] -> total=4 passed=3 failed=1 rate=75.0
\end{pygym}
\textbf{Note.} The aggregate behind every test report; rate is what goes at the
top of the summary.
\end{gymproblem}

%---------------------------------------------------------------
\begin{gymproblem}{Hexdump a byte buffer}
\begin{pygym}
def hexdump(b):
    return ' '.join('%02X' % x for x in b)
# b'\xDE\xAD\xBE\xEF' -> 'DE AD BE EF'
\end{pygym}
\textbf{Note.} The first tool you reach for when \code{unpack} gives nonsense ---
decode the raw bytes by hand to anchor the format (Q2.2).
\end{gymproblem}

%---------------------------------------------------------------
\begin{gymproblem}{Parse a CSV telemetry line}
\begin{pygym}
def parse_csv_line(line):
    return [float(x) for x in line.strip().split(',')]
# "1.5, 2.0, -3.25" -> [1.5, 2.0, -3.25]
\end{pygym}
\textbf{Note.} Text telemetry logs are often CSV; \code{strip()} tolerates
whitespace. For big files prefer the \code{csv} module.
\end{gymproblem}

%---------------------------------------------------------------
\begin{gymproblem}{Validate a frame checksum}
\begin{pygym}
def checksum_ok(frame):
    s = sum(frame[:-1]) & 0xFF
    return ((-s) & 0xFF) == frame[-1]   # last byte = two's-complement of the sum
\end{pygym}
\textbf{Note.} Mirrors the C \code{checksum} (Chapter 13); validate \emph{before}
trusting any field, and discard/resync on failure.
\end{gymproblem}

%---------------------------------------------------------------
\begin{gymproblem}{Compare with tolerance}
\begin{pygym}
def within_tol(expected, actual, tol):
    return abs(expected - actual) <= tol
# within_tol(100.0, 100.4, 0.5) -> True
\end{pygym}
\textbf{Note.} Hardware values are never bit-exact (ADC noise, timing) --- assert
a band, not equality, exactly as a voter does (Chapter 8/10).
\end{gymproblem}

%---------------------------------------------------------------
\begin{gymproblem}{Byte-swap a 32-bit value}
\begin{pygym}
def bswap32(v):
    return ((v & 0xFF) << 24) | ((v & 0xFF00) << 8) | \
           ((v >> 8) & 0xFF00) | ((v >> 24) & 0xFF)
# bswap32(0x11223344) -> 0x44332211
\end{pygym}
\textbf{Note.} When you must convert byte order outside \code{struct} (e.g.\ a
field you already have as an int); matches the C \code{bswap32}.
\end{gymproblem}

%---------------------------------------------------------------
\begin{gymproblem}{Write a CSV report}
\begin{pygym}
import io, csv

def make_csv(rows):
    buf = io.StringIO()
    w = csv.writer(buf)
    w.writerow(['name', 'expected', 'actual', 'pass'])
    for r in rows:
        w.writerow(r)
    return buf.getvalue()
\end{pygym}
\textbf{Note.} The \code{csv} module handles quoting/escaping correctly --- never
hand-format CSV. The output imports into a traceability tool or spreadsheet.
\end{gymproblem}

%---------------------------------------------------------------
\begin{gymproblem}{Generate a markdown report}
\begin{pygym}
def make_report(name, results):
    s = summarize(results)
    return f"# {name}\n\nPassed {s['passed']}/{s['total']} ({s['rate']}%)\n"
# -> "# UART Test\n\nPassed 2/3 (66.7%)\n"
\end{pygym}
\textbf{Note.} A human-readable artefact for the test record; f-strings keep it
readable. Extend with a per-test table.
\end{gymproblem}

%---------------------------------------------------------------
\begin{gymproblem}{Parse a key=value config line}
\begin{pygym}
def parse_kv(line):
    d = {}
    for pair in line.split():
        if '=' in pair:
            k, v = pair.split('=', 1)
            d[k] = v
    return d
# "baud=115200 parity=none" -> {'baud': '115200', 'parity': 'none'}
\end{pygym}
\textbf{Note.} Data-driven test configuration --- the maintainable alternative to
hard-coded parameters (Q2.5).
\end{gymproblem}

%---------------------------------------------------------------
\begin{gymproblem}{Serial command/response (\texttt{\# requires pyserial + target})}
\begin{pygym}
# requires: pip install pyserial ; and a connected device
import serial, struct

def query(port, fid, seq, value):
    with serial.Serial(port, 115200, timeout=1) as s:   # timeout: never block forever
        s.write(struct.pack('>BBf', fid, seq, value))   # send command frame
        resp = s.read(8)                                # read fixed-size response
        if len(resp) < 8:
            raise TimeoutError("no/short response")     # timeout = defined failure
        return struct.unpack('>BBfH', resp)
\end{pygym}
\textbf{Note.} The end-to-end automation step: pack $\rightarrow$ write
$\rightarrow$ read-with-timeout $\rightarrow$ unpack. The \code{with} block closes
the port even on exception; a short read is a test failure, not a hang (Q2.3).
\end{gymproblem}

%=====================================================================
\section{Link to Her Resume}
%=====================================================================
\begin{resumelink}
\textbf{The 60\% automation suite:} ``My Python suite parsed MIL-STD-1553B and
RS422 telemetry and cut manual regression effort ~60\%. The core is
\code{struct} pack/unpack with explicit big-endian framing, sub-byte field
extraction by shift/mask, CRC validation matched to the firmware's exact
parameters, and a data-driven test layer that turns each case into a table row
--- so the full suite runs unattended and repeatably.''

\medskip
\textbf{Built, not just run:} ``The hard parts were byte-level: framing and
resync on a noisy serial link, timeouts so the suite never hangs, validating the
status/SSM before trusting a value, and matching the CRC poly/init/reflection to
the target. I keep Python as the host-side controller/analyser and leave
hard-real-time timing to the hardware/HIL rig.''

\medskip
\small Maps directly to the resume's Python automation (1553B/RS422 parsing,
~60\% regression saving) and SAP QM reporting. Deep project scripting is
Chapter 15, from \code{inputs/INTAKE.md} only.
\end{resumelink}

%=====================================================================
\section{Self-Test Scorecard}
%=====================================================================
\begin{scorecard}
\begin{enumerate}
\item What module handles binary protocol data, and what does \code{>BBfH} mean?
\item Why is the endianness prefix mandatory, and what breaks without it?
\item How do you extract a 5-bit field that \code{struct} can't unpack directly?
\item Give the \code{struct.unpack}-gives-nonsense debug checklist.
\item Why must serial reads have a timeout?
\item Name the CRC parameters you must match between Python and firmware.
\item Why assert a tolerance band rather than equality on hardware values?
\item What makes an automation framework maintainable?
\item Where is Python the wrong tool in a hardware test, and what replaces it?
\item How do you \emph{prove} you built the automation, not just ran it?
\end{enumerate}
\end{scorecard}

\begin{answerkey}
\begin{enumerate}
\item \code{struct}; big-endian unsigned byte, byte, 32-bit float, 16-bit
  unsigned --- an 8-byte frame.
\item It sets byte order; without the explicit prefix you get native order
  \emph{and} padding, byte-swapping/misaligning every multi-byte field.
\item Unpack the containing integer with \code{struct}, then shift and mask
  (e.g.\ \code{(w>>11)\&0x1F}).
\item Endianness, format vs actual sizes/order, padding (use \code{>}/\code{<}),
  byte length vs \code{calcsize}, and true frame start (sync) --- hexdump to
  anchor.
\item So a test never blocks forever on missing hardware data; a timeout is a
  defined failure.
\item Polynomial, initial value, input/output reflection, final XOR, and the
  covered byte range.
\item ADC noise/timing make values non-exact; equality nuisance-fails, a band
  reflects real tolerance.
\item Layered design (protocol/transport/tests/reporting), data-driven cases, a
  real test runner, versioned expected results.
\item For hard-real-time/cycle-accurate timing (GC/interpreter jitter); put that
  in hardware/firmware/HIL with Python as controller.
\item Describe the design choices and edge cases --- framing/resync, CRC-parameter
  matching, timeouts/partial frames, SSM validation, maintainability --- and the
  measured outcome.
\end{enumerate}
\end{answerkey}

\begin{partnote}
\emph{End of Chapter 14, and of Part D (Coding Gym).} Part E (Interview
Execution) follows. Chapters 15--16 (Her Project Deep-Dives; Behavioural \&
Company) are written \textbf{only} from \code{inputs/INTAKE.md} and are gated
until it is complete --- so the next written chapter is \textbf{Chapter 17: Mock
Interviews}.
\end{partnote}
